problem:  
Let \( (M^7, \varphi) \) be a compact Riemannian manifold with torsion-free \( G_2 \)-structure and a **free, isometric circle action** generated by a Killing vector field \(V\) which preserves \(\varphi\). Denote by \(\pi: M \to X = M/S^1\) the quotient 6-manifold and by \(\theta\) a connection 1-form on the principal \(S^1\)-bundle \(M \to X\) with curvature \(F = d\theta\).  

It is known (Apostolov–Salamon) that under these assumptions, the pair \((\omega, \Omega)\) of differential forms on \(X\), defined by descending
\[
\omega = \iota_V \varphi, \quad \text{and} \quad \Re \Omega = \varphi|_X, \quad \Im \Omega = -\, \iota_V {*\varphi},
\]
satisfies a nonlinear system of PDEs that ensure \(\varphi\) is torsion-free.  

**Conjectural Problem:**  
Is every compact torsion-free \(G_2\)-manifold \( (M, \varphi) \) admitting such a free isometric circle symmetry obtained as the total space of a principal \(S^1\)-bundle over a compact Calabi–Yau 3-fold \( (X, \omega, \Omega) \) satisfying the Apostolov–Salamon equations?  
More precisely, does there exist a universal “reconstruction” correspondence 
\[
\{\text{Calabi–Yau 3-folds } X \text{ with certain } (F, \theta)\text{ satisfying Apostolov–Salamon equations}\}
\;\longleftrightarrow\;
\{\text{compact torsion-free } G_2 \text{ manifolds with free } S^1 \text{ symmetry}\},
\]
up to diffeomorphism and scaling of \(\varphi\)?  

In other words, can all such \(G_2\)-manifolds be derived from Calabi–Yau geometry via the Apostolov–Salamon reduction ansatz, possibly after connected-sum or generalized Kummer-type gluing?  

Why is it a "good" problem:  
This problem sits precisely at the interface of **exceptional geometry** and **complex/symplectic geometry**, probing whether all \( G_2 \)-manifolds with symmetry stem from Calabi–Yau sources via a fully geometric reduction mechanism. The question is conceptually simple—it asks whether “every symmetric \( G_2 \)-structure hides a Calabi–Yau base”—yet it opens into deep analytic, topological, and geometric territory.  

Resolving it would potentially classify an entire class of compact \( G_2 \)-manifolds by a lower-dimensional complex-geometric moduli problem, offering a bridge between special holonomy, gauge theory, and complex geometry. It would extend known partial constructions (via Apostolov–Salamon, twisted connected sums, or Joyce’s desingularization) into a **unified geometric correspondence**.  
Thus, the problem embodies the essence of a “good” research problem: technically sound, simple to state, grounded in existing deep structures, but pointing to an unknown and potentially unifying principle.